%% kapa: exam=1 round=28 year=2017 period=2 form=A subject=accounting subject_ko=회계학 seq=11 number=51
%% source: Q-Net 감정평가사 기출문제 (cst00309) artlSeq=5199603 file=extracted/1008/5199603/2017년 제28회 감정평가사 1차시험 문제지/2교시 A형.pdf
%% conversion: kordoc 4.24.51 -> tools/gen1.py (auto); review=proofread-sonnet; flags=none
\begin{question}{51}
(주)감평은 20x1년 초에 도급금액 ￦1,000,000인 건설공사를 수주하고, 20x3년 말에 공사를 완공하였다. 이와 관련된 원가자료는 다음과 같다. (주)감평이 20x1년도 포괄손익계산서에 인식할 공사손익과 20x1년말 재무상태표에 표시할 미청구공사(또는 초과청구공사) 금액은? (단, 진행률은 발생누적계약원가를 추정총계약원가로 나눈 비율로 계산한다.)
\par{\small\centering\begin{tabular}{|l|c|c|c|}
\hline
 & 20x1년 & 20x2년 & 20x3년 \\ \hline
실제발생 공사원가 & ￦320,000 & ￦200,000 & ￦250,000 \\ \hline
연도 말 예상 추가원가 & 480,000 & 280,000 & - \\ \hline
계약대금 청구액 & 350,000 & 350,000 & 300,000 \\ \hline
\end{tabular}\par}
\begin{center}
\begin{tabular}{@{}lll@{}}
 & \uline{공사이익(손실)} & \uline{미청구공사(초과청구공사)} \\
① & ￦80,000 & ￦ 50,000 \\
② & ￦60,000 & ￦ 30,000 \\
③ & ￦60,000 & ￦(30,000) \\
④ & ￦80,000 & ￦(50,000) \\
⑤ & ￦80,000 & ￦ 30,000 \\
\end{tabular}
\end{center}
\end{question}
\answer{1}
